\documentclass[../DoAn.tex]{subfiles}
\begin{document}

Chương này tập trung vào những giải pháp mà tác giả đánh giá là có giá trị nhất trong suốt quá trình thực hiện đồ án. Khác với Chương 4 vốn đặt trọng tâm vào mô tả thiết kế, hiện thực và kết quả của toàn hệ thống, chương này đi sâu vào các bài toán khó, các quyết định kỹ thuật quan trọng và các cách tiếp cận đã được lựa chọn để giải quyết các bài toán đó. Mỗi mục dưới đây không nhằm liệt kê lại chức năng đã có, mà nhằm làm rõ vì sao bài toán đó khó, giải pháp nào đã được áp dụng, và giải pháp đó đã mang lại giá trị gì cho sản phẩm.

Trong phạm vi đồ án này, bốn đóng góp nổi bật nhất gồm: giải pháp thích ứng đa nền tảng cho các chức năng phụ thuộc thiết bị; giải pháp bảo toàn tính nhất quán cho luồng đặt vé nhiều hạng kết hợp thanh toán và check-in bằng mã \gls{QR}; giải pháp xây dựng cơ chế nhắc lịch và thông báo hợp nhất, hạn chế gửi trùng; và giải pháp tích hợp mô-đun AI hỗ trợ tạo nội dung sự kiện với cơ chế dự phòng nhiều tầng. Đây đều là những vấn đề có liên quan trực tiếp tới khả năng sử dụng thật của hệ thống, đồng thời thể hiện rõ nhất tư duy phân tích và tổng quát hóa trong quá trình làm đồ án.

\section{Giải pháp thích ứng đa nền tảng cho các chức năng phụ thuộc thiết bị}
\label{sec:solution-cross-platform}

\subsection{Bài toán và bối cảnh}
Một trong những mục tiêu thực tế của dự án là không chỉ dừng ở ứng dụng mobile chạy trên Android/iOS, mà còn có thể khai thác lại mã nguồn để phục vụ bản chạy trên web. Tuy nhiên, khi chuyển một ứng dụng React Native sang môi trường web, những chức năng phụ thuộc chặt vào thiết bị thường là nơi phát sinh nhiều vấn đề nhất. Trong hệ thống event-booking-app, hai nhóm chức năng nhạy cảm nhất là chọn vị trí trên bản đồ và quét vé bằng camera. Trên thiết bị di động, ứng dụng có thể dùng trực tiếp \texttt{react-native-maps}, \texttt{expo-location} hoặc camera của Expo. Trên web, các thư viện này không còn giữ nguyên hành vi như trên mobile, hoặc chịu thêm ràng buộc của trình duyệt như \gls{HTTPS}, quyền truy cập camera và cơ chế render khác biệt.

Nếu giải quyết theo cách đơn giản nhất, người phát triển có thể bỏ hẳn các chức năng này khỏi web, hoặc chỉ hỗ trợ ở mức rất hạn chế. Tuy nhiên, cách làm đó khiến giá trị đa nền tảng của hệ thống bị giảm mạnh, đặc biệt trong bối cảnh web lại là môi trường thuận tiện để trình diễn sản phẩm hoặc triển khai thử nghiệm khi chưa đóng gói ứng dụng di động hoàn chỉnh. Ngược lại, nếu viết lại toàn bộ một phiên bản web tách biệt với mobile thì chi phí bảo trì cao, luồng nghiệp vụ dễ sai khác, và bản thân đồ án sẽ mất đi lợi thế “một lõi nghiệp vụ, nhiều điểm truy cập”. Do đó, bài toán cốt lõi không chỉ là làm cho web chạy được, mà là làm sao thích ứng với khác biệt nền tảng mà vẫn giữ được sự thống nhất ở luồng sử dụng và dữ liệu trao đổi với backend.

\subsection{Giải pháp}
Giải pháp được lựa chọn là tách phần phụ thuộc nền tảng ra khỏi phần nghiệp vụ dùng chung, tức là giữ nguyên các service, route điều hướng, dữ liệu đầu vào/đầu ra và quy tắc nghiệp vụ, nhưng thay thế lớp tương tác thiết bị bằng các adapter riêng cho từng môi trường. Cách tiếp cận này được thể hiện khá rõ qua hai cặp màn hình \texttt{SelectLocation.tsx} và \texttt{SelectLocation.web.tsx}, cũng như \texttt{CheckIn.tsx} và \texttt{CheckIn.web.tsx}. Ở đây, tác giả không tách hai ứng dụng độc lập, mà chỉ tách lớp hiện thực giao diện và truy cập thiết bị, còn các tham số điều hướng, cách gọi service, cấu trúc dữ liệu và trạng thái luồng nghiệp vụ vẫn được giữ tương đồng.

Đối với bài toán chọn vị trí, bản mobile sử dụng \texttt{MapView} kết hợp \texttt{expo-location} để lấy tọa độ hiện tại, di chuyển camera và xác nhận vị trí trên bản đồ. Trong khi đó, bản web chuyển sang dùng Leaflet và cơ chế \texttt{navigator.geolocation}, đồng thời nạp động stylesheet và icon tương ứng cho môi trường trình duyệt. Điểm quan trọng không nằm ở việc thay thư viện bản đồ, mà ở chỗ toàn bộ quy trình nghiệp vụ vẫn được giữ nguyên: người dùng có thể tìm kiếm địa chỉ qua Nominatim, chọn kết quả gợi ý, tinh chỉnh vị trí trên bản đồ, reverse geocode để lấy tên địa điểm, rồi trả kết quả về đúng màn hình gọi trước đó. Như vậy, sự khác biệt chỉ nằm ở “cách người dùng thao tác với bản đồ”, còn “ý nghĩa nghiệp vụ của bước chọn vị trí” không thay đổi.

Tương tự, với chức năng check-in vé, bản mobile ưu tiên quét camera trực tiếp thông qua \texttt{expo-camera}. Tuy nhiên, trên web, camera thường gặp các ràng buộc như không chạy trên \gls{HTTP}, bị người dùng từ chối cấp quyền, hoặc hoạt động không ổn định giữa các trình duyệt. Thay vì xem đây là lý do để bỏ chức năng, hệ thống bổ sung một nhánh dự phòng cho web là nhập mã \gls{QR} thủ công. Cách làm này cho thấy tư duy thiết kế tập trung vào nghiệp vụ thay vì phụ thuộc tuyệt đối vào thiết bị: bản chất của check-in không phải là “phải quét bằng camera”, mà là “xác minh dữ liệu vé và cập nhật trạng thái tham dự một cách an toàn”. Camera chỉ là một phương tiện thuận tiện để đưa dữ liệu vào hệ thống.

Điểm quan trọng hơn là dù giao diện nền tảng khác nhau, cả mobile lẫn web đều gọi chung các service nghiệp vụ và cùng đi tới một backend duy nhất. Điều đó giúp tránh tình trạng mỗi nền tảng tự phát triển một logic kiểm tra riêng. Khi thay đổi quy tắc xác minh check-in hoặc quy tắc lấy vị trí, tác giả chỉ cần cập nhật một tầng nghiệp vụ thay vì phải đồng bộ thủ công ở nhiều nơi. Có thể xem đây là một giải pháp thích ứng theo lớp: tầng giao diện chấp nhận khác biệt nền tảng, còn tầng service và tầng backend vẫn giữ vai trò nguồn chân lý duy nhất của hệ thống.

\subsection{Kết quả đạt được}
Kết quả rõ nhất của giải pháp này là hệ thống duy trì được khả năng chạy đa nền tảng mà không phải hy sinh các chức năng quan trọng. Người dùng mobile vẫn có trải nghiệm tự nhiên với bản đồ và camera thiết bị, trong khi người dùng web vẫn sử dụng được các chức năng tương đương với một số điều chỉnh hợp lý theo đặc thù trình duyệt. Điều này làm tăng đáng kể giá trị thực tế của sản phẩm trong giai đoạn thử nghiệm và trình diễn.

Về mặt kỹ thuật, giải pháp giúp giảm trùng lặp mã nguồn nghiệp vụ. Những gì thực sự khác biệt mới bị tách riêng sang tệp \texttt{.web.tsx}; các phần như điều hướng, service, xác thực, gọi \gls{API} và cấu trúc dữ liệu vẫn được tái sử dụng. Nhờ đó, chi phí bảo trì thấp hơn so với việc viết hai phiên bản độc lập. Quan trọng hơn, giải pháp này có thể tổng quát hóa cho nhiều chức năng khác của ứng dụng đa nền tảng: khi gặp khác biệt phần cứng hoặc môi trường, không nhất thiết phải bỏ chức năng hoặc nhân đôi toàn bộ hệ thống, mà có thể giữ một lõi nghiệp vụ chung và chỉ thay adapter tương tác ở lớp biên.

\section{Giải pháp bảo toàn tính nhất quán cho luồng đặt vé nhiều hạng kết hợp thanh toán và check-in}
\label{sec:solution-booking-consistency}

\subsection{Bài toán và bối cảnh}
Nghiệp vụ trung tâm của hệ thống là cho phép người dùng đặt vé tham gia sự kiện. Nếu chỉ dừng ở mức mỗi sự kiện có một mức giá cố định và không có xác nhận tham dự, bài toán tương đối đơn giản. Tuy nhiên, hệ thống của đồ án đi theo hướng thực tế hơn: mỗi sự kiện có thể có nhiều hạng vé với giá và quota khác nhau; người dùng có thể đặt nhiều vé trong một giao dịch; sau thanh toán, hệ thống phải phát hành vé điện tử; khi đến sự kiện, ban tổ chức phải kiểm tra được vé nào hợp lệ, vé nào đã dùng, vé nào thuộc sự kiện khác hoặc chưa thanh toán. Chính sự nối tiếp của nhiều bước này khiến bài toán không còn là “lưu một bản ghi mua vé”, mà là bài toán bảo toàn tính nhất quán của cả một chuỗi nghiệp vụ.

Khó khăn lớn nhất nằm ở việc dữ liệu từ phía client không thể được tin cậy hoàn toàn. Nếu frontend gửi lên mức giá đã tính sẵn, backend chấp nhận trực tiếp thì rất dễ phát sinh nguy cơ sửa giá ở phía client. Nếu vé được sinh ngay khi người dùng bấm đặt mà chưa gắn với trạng thái thanh toán, hệ thống có thể phát hành vé chưa trả tiền. Nếu check-in chỉ kiểm tra mã tồn tại mà không kiểm tra sự kiện, trạng thái thanh toán hoặc trạng thái đã dùng, người tổ chức có thể quét nhầm hoặc chấp nhận vé không hợp lệ. Như vậy, thách thức ở đây không nằm ở một hàm đơn lẻ, mà nằm ở việc thiết kế luồng dữ liệu sao cho mỗi bước đều kiểm tra lại các giả định quan trọng.

\subsection{Giải pháp}
Giải pháp được áp dụng là dồn toàn bộ quy tắc kiểm tra quan trọng về backend và tách nghiệp vụ thành các pha rõ ràng: kiểm tra dữ liệu đặt vé, tạo vé ở trạng thái chờ, tạo giao dịch thanh toán, xác nhận thanh toán, sau đó mới phát hành dữ liệu \gls{QR} và cho phép check-in. Trong controller đặt vé, backend không sử dụng giá client gửi lên như một nguồn chân lý tuyệt đối. Thay vào đó, hệ thống tra cứu hạng vé theo \texttt{tierName} ngay trên tài liệu sự kiện, tính lại \texttt{expectedTotal} từ giá của hạng vé và số lượng vé, rồi so sánh với số tiền được gửi từ client. Nếu hai giá trị không khớp, yêu cầu bị từ chối. Đây là một bước nhỏ nhưng rất quan trọng, vì nó chặn được lớp lỗi phổ biến nhất ở các ứng dụng thương mại điện tử đơn giản: tin vào giá phía client.

Sau khi xác thực giá và quota, hệ thống tạo nhiều bản ghi \texttt{Ticket} ứng với số lượng vé mua, nhưng tất cả đều bắt đầu ở trạng thái \texttt{paymentStatus = pending}. Đồng thời, hệ thống chỉ tạo một bản ghi \texttt{Payment} duy nhất cho toàn giao dịch. Cách tổ chức này mang lại hai lợi ích. Thứ nhất, người dùng vẫn có thể mua nhiều vé trong một lần thao tác. Thứ hai, hệ thống phân biệt được “đơn vị quyền tham dự” là từng vé, còn “đơn vị giao dịch tài chính” là một thanh toán. Đây là tách biệt có ý nghĩa dài hạn, bởi nếu sau này cần hoàn tiền, thống kê doanh thu hoặc tích hợp cổng thanh toán thật, lớp \texttt{Payment} có thể được mở rộng mà không phá vỡ mô hình vé.

Khi thanh toán được xác nhận thành công, backend mới cập nhật hàng loạt trạng thái vé sang \texttt{paid}, tăng số lượng vé đã bán của đúng hạng vé tương ứng, tính lại tổng số người tham gia của sự kiện, tạo thông báo thanh toán và sinh dữ liệu \gls{QR} cho từng vé. Việc sinh \gls{QR} ở giai đoạn sau thanh toán thay vì ở giai đoạn đặt giữ cho hệ thống chặt chẽ hơn: chỉ vé đã thanh toán mới có tư cách trở thành vé điện tử hợp lệ. Dữ liệu \gls{QR} được thiết kế gọn theo dạng \texttt{ticketId|eventId|token}, nghĩa là mã quét không chỉ trỏ tới một vé, mà còn gắn chặt với một sự kiện cụ thể và có thêm token phụ trợ để giảm tính đoán trước.

Phần check-in tiếp tục tuân theo nguyên tắc “mọi xác nhận quan trọng đều thực hiện ở backend”. Khi người tổ chức quét hoặc nhập dữ liệu mã, backend lần lượt kiểm tra: người thực hiện có đúng là organizer của sự kiện hay không; mã có đúng định dạng không; vé trong mã có thuộc đúng sự kiện đang check-in không; vé có thực sự tồn tại không; vé đã thanh toán chưa; vé đã check-in trước đó chưa. Chỉ khi vượt qua toàn bộ các bước này, hệ thống mới cập nhật trạng thái \texttt{checked}, thời điểm \texttt{checkedAt} và người thực hiện \texttt{checkedBy}. Nhờ vậy, check-in không chỉ là thao tác đánh dấu, mà là bước xác minh cuối cùng của toàn chuỗi từ đặt vé tới tham dự.

\subsection{Kết quả đạt được}
Giải pháp này giúp hệ thống đạt được một mức nhất quán nghiệp vụ khá tốt trong phạm vi đồ án. Các nguy cơ như sửa giá ở phía client, đặt vượt quota của hạng vé, sinh vé trước khi thanh toán, dùng vé sai sự kiện hoặc check-in lặp lại đều đã được chặn bằng các bước xác thực cụ thể ở phía backend. So với cách cài đặt “frontend tính gì backend lưu nấy”, cách tiếp cận hiện tại an toàn hơn và có tính mở rộng tốt hơn.

Một kết quả khác là luồng nghiệp vụ trở nên rõ ràng và dễ kiểm thử hơn. Vì mỗi pha đều có trách nhiệm riêng, việc viết use case, biểu đồ trình tự và test cho các nhánh thành công hoặc thất bại cũng thuận lợi hơn. Ở góc độ tổng quát, đây là một giải pháp có thể áp dụng cho nhiều hệ thống đặt chỗ hoặc bán vé khác: không để nghiệp vụ quan trọng trộn lẫn trong một bước duy nhất, mà chia thành các pha có trạng thái trung gian rõ ràng, để mỗi bước đều có thể được xác minh và phục hồi khi cần.

\section{Giải pháp xây dựng cơ chế nhắc lịch và thông báo hợp nhất, hạn chế gửi trùng}
\label{sec:solution-reminder-notification}

\subsection{Bài toán và bối cảnh}
Đối với một ứng dụng sự kiện, trải nghiệm của người dùng không kết thúc ở thời điểm đặt vé. Sau khi đã quan tâm hoặc mua vé, người dùng còn cần được nhắc nhớ đúng lúc để không bỏ lỡ sự kiện; người tổ chức cũng cần nhận tín hiệu về các hành vi quan trọng như thanh toán thành công hoặc tương tác từ người dùng. Nếu không có lớp thông báo phù hợp, ứng dụng rất dễ trở thành một nơi “đặt xong rồi quên”, khiến giá trị sử dụng thực tế giảm mạnh.

Vấn đề là thông báo trong hệ thống không phải chỉ có một loại. Ở dự án này, có thông báo thanh toán, lời mời sự kiện, nhắc lịch trước giờ diễn ra, và các cập nhật liên quan tới chat. Nếu mỗi loại được xử lý theo một kiểu khác nhau, hệ thống sẽ nhanh chóng trở nên rời rạc: nơi lưu ở database, nơi chỉ gửi push, nơi không kiểm soát việc gửi trùng, nơi không quan tâm người dùng đã tắt thông báo hay chưa. Đặc biệt với thông báo nhắc lịch, nguy cơ gửi trùng là khá lớn, nhất là khi scheduler chạy định kỳ hoặc khi server khởi động lại.

\subsection{Giải pháp}
Giải pháp được lựa chọn là xây dựng một trục thông báo hợp nhất, trong đó \texttt{Notification} đóng vai trò thực thể trung tâm để lưu toàn bộ thông báo nội bộ, còn push notification là lớp phát bổ sung ra thiết bị khi có token phù hợp. Với cách làm này, mọi thông báo quan trọng đều có dấu vết trong hệ thống, giúp người dùng có thể xem lại trong giao diện thông báo kể cả khi push bị trượt hoặc không mở được ngay thời điểm nhận.

Riêng đối với nhắc lịch sự kiện, tác giả triển khai một scheduler chạy theo chu kỳ 5 phút. Thay vì chỉ lọc sự kiện theo ngày đơn thuần, hệ thống tách riêng \texttt{date} và \texttt{time} rồi ghép lại thành thời điểm bắt đầu thực tế của sự kiện. Sau đó, scheduler xét hai cửa sổ thời gian quan trọng là trước 1 ngày và trước 1 giờ. Cách tiếp cận này giải quyết một vấn đề thường gặp khi lưu ngày và giờ tách rời: nếu chỉ so sánh theo trường ngày, hệ thống khó xác định chính xác thời điểm nhắc cho những sự kiện diễn ra vào các khung giờ khác nhau trong cùng một ngày.

Để tránh gửi trùng, giải pháp sử dụng hai lớp bảo vệ. Lớp thứ nhất là một \texttt{Map} trong bộ nhớ để ghi nhận những thông báo đã được gửi ở cửa sổ hiện tại, tránh việc scheduler lặp lại cùng một hành động trong thời gian ngắn. Tuy nhiên, chỉ dựa vào bộ nhớ là chưa đủ vì khi server khởi động lại thì trạng thái này mất đi. Do đó, lớp thứ hai là truy vấn ngược vào collection \texttt{Notification} để kiểm tra xem với cùng người dùng, cùng sự kiện và cùng nội dung nhắc, hệ thống đã tạo thông báo hay chưa. Chỉ những người nhận chưa có thông báo tương ứng mới tiếp tục được gửi push và ghi mới vào database. Đây là một điểm mà tác giả đánh giá là quan trọng, vì nó biến cơ chế nhắc lịch từ “gần đúng” thành “đủ tin cậy để dùng thật”.

Một quyết định thiết kế khác là xác định người nhận theo nghĩa nghiệp vụ thay vì chỉ dựa trên danh sách nhắc thủ công. Với mỗi sự kiện sắp diễn ra, hệ thống lấy hợp của hai tập: người tổ chức sự kiện và những người đã có vé ở trạng thái \texttt{paid}. Cách xác định này hợp lý hơn cho luồng hiện tại vì người thực sự cần được nhắc là người có trách nhiệm tổ chức và người chắc chắn sẽ tham dự. Trên cùng nền tảng đó, lời mời sự kiện và thông báo thanh toán cũng được ghi vào chung hệ thống thông báo, từ đó giao diện notification ở frontend có thể hỗ trợ đồng nhất các thao tác như xem danh sách, đánh dấu đã đọc, đánh dấu một phần hoặc xóa thông báo.

\subsection{Kết quả đạt được}
Kết quả của giải pháp là hệ thống có được một lớp thông báo tương đối thống nhất giữa backend và frontend, đồng thời tạo nền tảng tốt để mở rộng sau này. Người dùng không chỉ nhận push notification khi có điều kiện phù hợp, mà còn có một hộp thư thông báo nội bộ để tra cứu lại các sự kiện đã xảy ra. Điều này đặc biệt hữu ích với lời mời tham gia sự kiện hoặc nhắc lịch trước giờ diễn ra.

Ở góc độ kỹ thuật, giá trị lớn nhất của giải pháp là hạn chế được hiện tượng gửi trùng vốn rất dễ xuất hiện trong các hệ thống scheduler đơn giản. Việc kết hợp giữa bộ nhớ tạm và kiểm tra database cho thấy một hướng xử lý thực tế: không kỳ vọng một cơ chế đơn lẻ giải quyết toàn bộ bài toán, mà kết hợp nhiều lớp kiểm soát với chi phí hợp lý. Đây là một kinh nghiệm có thể tổng quát hóa cho các bài toán gửi thông báo định kỳ khác, không chỉ trong ứng dụng sự kiện mà còn trong các ứng dụng nhắc việc, lớp học, đặt lịch hay chăm sóc khách hàng.

\section{Giải pháp tích hợp AI hỗ trợ tạo nội dung sự kiện với cơ chế dự phòng nhiều tầng}
\label{sec:solution-ai-fallback}

\subsection{Bài toán và bối cảnh}
Một khó khăn thực tế của người tổ chức sự kiện, đặc biệt ở những ứng dụng hướng cộng đồng, là họ thường không gặp khó trong việc chọn ngày hay địa điểm, mà gặp khó ở khâu mô tả nội dung sao cho rõ ràng, hấp dẫn và nhất quán. Nếu để trống hoàn toàn bước này, nhiều sự kiện dễ có tiêu đề sơ sài, mô tả nghèo nàn hoặc hạng vé được đặt tên thiếu hợp lý, từ đó làm giảm chất lượng chung của nền tảng. Vì vậy, tác giả mong muốn hệ thống không chỉ là nơi “đăng sự kiện”, mà còn hỗ trợ người tổ chức hình thành nội dung ban đầu.

Tuy nhiên, khi đưa AI vào một đồ án theo hướng ứng dụng, bài toán không dừng ở việc gọi được một mô hình sinh văn bản. Dịch vụ AI bên ngoài có thể lỗi, hết hạn khóa, trả về dữ liệu sai cấu trúc hoặc có độ trễ lớn. Nếu frontend phụ thuộc hoàn toàn vào một lời gọi không ổn định như vậy, trải nghiệm tạo sự kiện sẽ bị ảnh hưởng mạnh. Nói cách khác, thách thức không chỉ là “có AI”, mà là “tích hợp AI theo cách không làm hỏng luồng nghiệp vụ chính”.

\subsection{Giải pháp}
Giải pháp được chọn là đưa AI về phía backend dưới dạng một service độc lập, đồng thời áp dụng cơ chế fallback nhiều tầng. Ở tầng đầu tiên, hệ thống có thể dùng OpenAI làm nhà cung cấp ưu tiên. Nếu cấu hình yêu cầu, hệ thống cũng hỗ trợ gọi Ollama như một mô hình nội bộ hoặc cục bộ. Nếu cả hai cách đều không khả dụng, hệ thống vẫn không trả lỗi ngay lập tức mà rơi về một tầng cuối là sinh gợi ý theo template. Như vậy, luồng tạo sự kiện không bao giờ bị chặn hoàn toàn chỉ vì một nhà cung cấp AI không phản hồi.

Một điểm quan trọng khác của giải pháp là không nhận nội dung AI ở dạng văn bản tự do, mà ép đầu ra về một cấu trúc \gls{JSON} có schema rõ ràng gồm ba thành phần: tiêu đề, mô tả và danh sách hạng vé. Sau khi nhận dữ liệu, backend còn tiếp tục chuẩn hóa và kiểm tra lại để bảo đảm tên hạng vé không rỗng, giá là số hợp lệ, quota dương và số lượng hạng vé nằm trong phạm vi cho phép. Cách làm này giúp giảm đáng kể rủi ro một mô hình sinh ra câu trả lời đẹp nhưng không dùng được trực tiếp trong hệ thống. Về bản chất, AI chỉ đóng vai trò đề xuất; quyền kiểm soát cấu trúc dữ liệu vẫn thuộc về lớp nghiệp vụ.

Ở phía frontend, mô-đun này được tích hợp ngay trong màn hình tạo sự kiện. Người tổ chức có thể cung cấp danh mục sự kiện, vị trí, ngày dự kiến, nội dung mô tả hiện tại và một lời nhắc ngắn về định hướng chương trình. Sau khi nhận gợi ý, frontend tự động điền dữ liệu vào biểu mẫu và cho phép người dùng tiếp tục chỉnh sửa trước khi gửi tạo sự kiện thật. Điều này cho thấy AI không thay thế người tổ chức, mà hoạt động như một công cụ hỗ trợ tăng tốc bước khởi tạo nội dung.

Điểm mà tác giả đánh giá là có giá trị trong giải pháp này là tư duy “AI là một dịch vụ không hoàn toàn tin cậy”. Vì thế, toàn bộ cách tích hợp đều đi theo hướng chống phụ thuộc tuyệt đối: có schema để kiểm soát cấu trúc, có sanitization để kiểm soát dữ liệu, có fallback để giữ luồng nghiệp vụ không bị gián đoạn, và có khả năng chuyển đổi nhà cung cấp qua biến môi trường. Đây là hướng tiếp cận phù hợp hơn với ứng dụng thực tế so với việc gọi trực tiếp một API AI từ frontend và hy vọng mọi phản hồi đều hợp lệ.

\subsection{Kết quả đạt được}
Kết quả trước hết là người tổ chức có thêm một công cụ hỗ trợ rất thực dụng khi tạo sự kiện. Thay vì bắt đầu từ biểu mẫu trống, họ có thể nhận được một bộ khung ban đầu gồm tiêu đề, mô tả và 1--3 hạng vé hợp lý để tiếp tục tinh chỉnh. Điều này giúp giảm đáng kể ma sát ở bước khởi tạo, nhất là với các sự kiện cộng đồng nhỏ hoặc các trường hợp người dùng chưa quen cách viết nội dung quảng bá.

Ở góc độ kỹ thuật, giải pháp chứng minh rằng có thể tích hợp AI vào một hệ thống nghiệp vụ mà không làm mất tính ổn định của hệ thống. Ngay cả khi không có khóa OpenAI, khi Ollama không hoạt động hoặc khi mạng có vấn đề, chức năng vẫn đưa ra được một gợi ý tối thiểu nhờ tầng template fallback. Đây là một đóng góp mà tác giả đánh giá là đáng giá vì nó chuyển AI từ vai trò “tính năng trình diễn” sang vai trò “thành phần có kiểm soát trong kiến trúc phần mềm”.

\section{Nhận xét chung về các đóng góp nổi bật}
Bốn giải pháp trên phản ánh bốn lát cắt khác nhau của cùng một bài toán xây dựng ứng dụng sự kiện. Giải pháp thứ nhất xử lý khác biệt nền tảng ở lớp biên; giải pháp thứ hai bảo đảm tính đúng đắn của luồng nghiệp vụ cốt lõi; giải pháp thứ ba tăng khả năng gắn kết và quay lại của người dùng thông qua thông báo; còn giải pháp thứ tư mở rộng hệ thống theo hướng thông minh nhưng vẫn giữ được tính ổn định. Nếu chỉ nhìn riêng từng tính năng, các kết quả này có thể xem là các chức năng đơn lẻ. Nhưng khi nhìn dưới góc độ kỹ thuật, chúng cho thấy một định hướng nhất quán: mọi quyết định đều cố gắng đưa hệ thống tiến gần hơn tới mức có thể sử dụng thật, thay vì chỉ dừng ở mức minh họa ý tưởng.

Từ quá trình thực hiện các giải pháp này, tác giả rút ra một nhận xét quan trọng: ở đồ án theo hướng sản phẩm, giá trị không nằm ở chỗ dùng bao nhiêu công nghệ mới, mà nằm ở chỗ xử lý được bao nhiêu tình huống “xấu” của bài toán thực tế. Khả năng chạy được trên nhiều môi trường, khả năng chống sai lệch dữ liệu từ client, khả năng tránh gửi thông báo trùng, hay khả năng giữ cho tính năng AI không phá vỡ luồng nghiệp vụ chính, đều là những ví dụ cho tư duy đó. Đây cũng chính là nhóm đóng góp mà tác giả xem là nổi bật nhất của đồ án.

\end{document}
